% Part: turing-machines
% Chapter: undecidability
% Section: introduction

\documentclass[../../../include/open-logic-section]{subfiles}

\begin{document}

\olfileid{tur}{und}{int}
\olsection{مقدمه}

شاید بدیهی به نظر برسد که هر تابعی، حتی هر تابع حسابی‌ای، نمی‌تواند
محاسبه‌پذیر باشد. شمار توابع بیش از اندازه زیاد و رفتارشان بیش از
اندازه پیچیده است؛ مثلاً توابعی که از واپاشی ذرات پرتوزا یا دیگر
رفتارهای آشوبناک یا تصادفی تعریف می‌شوند. فرض کنید از لحظه‌ای معین
شروع به شمردن بازه‌های یک‌ثانیه‌ای کنیم و $f(n)$ را تعداد ذرات جهان
تعریف کنیم که در $n$اُمین بازهٔ یک‌ثانیه‌ای پس از آن لحظهٔ آغازین
واپاشیده می‌شوند. به نظر می‌رسد این تابع نامزدی باشد که هرگز امیدی به
محاسبه‌اش نداریم.

اما نتوانستنِ تصور اینکه چنین توابعی چگونه محاسبه می‌شوند یک چیز است
و اثبات واقعیِ محاسبه‌ناپذیری‌شان چیزی دیگر. در واقع حتی توابعی که
بی‌اندازه پیچیده به نظر می‌رسند ممکن است در معنایی انتزاعی
محاسبه‌پذیر باشند. برای نمونه، فرض کنید جهان از نظر زمانی متناهی است:
روزی در آینده‌ای بسیار دور، آن‌گونه که برخی نظریه‌های کیهان‌شناختی
پیش‌بینی می‌کنند، جهان در یک نقطه منقبض خواهد شد. آنگاه از آن لحظهٔ
آغازین تنها تعداد متناهی (هرچند باورنکردنی بزرگی) ثانیه وجود دارد که
$f(n)$ برای آن‌ها تعریف شده است. هر تابعی که فقط برای تعداد متناهی
ورودی تعریف شده باشد محاسبه‌پذیر است: می‌توانیم خروجی‌ها را در جدولی
بزرگ فهرست کنیم یا آن‌ها را در نمودار گذار حالتِ ماشین تورینگی بسیار
بزرگ رمزگذاری کنیم.

اغلب به موارد ویژه‌ای از توابع علاقه‌مندیم که مقادیرشان پاسخ
پرسش‌های بله/خیر را می‌دهند. برای نمونه، پرسش «آیا $n$ عدد اول است؟»
با تابع زیر متناظر است:
\[
\fn{isprime}(n)=\begin{cases}
  1 & \text{اگر $n$ اول باشد}\\
  0 & \text{در غیر این صورت.}
  \end{cases}
\]
می‌گوییم یک پرسش بله/خیر \emph{به‌طور مؤثر تصمیم‌پذیر} است اگر تابع
متناظرِ $1/0$-مقداری به‌طور مؤثر محاسبه‌پذیر باشد.

برای آنکه به‌طور ریاضی ثابت کنیم توابعی هستند که نمی‌توان به‌طور مؤثر
محاسبه‌شان کرد یا مسائلی هستند که نمی‌توان به‌طور مؤثر درباره‌شان
تصمیم گرفت، ضروری است مدل مشخصی از محاسبه را تثبیت کنیم و نشان دهیم
توابع یا مسائلی وجود دارند که آن مدل نمی‌تواند محاسبه یا تصمیمشان
کند. برای نمونه، می‌توانیم نشان دهیم ماشین‌های تورینگ هر تابعی را
محاسبه و هر مسئله‌ای را تصمیم نمی‌کنند. سپس می‌توانیم با استناد به
تز چرچ--تورینگ نتیجه بگیریم که نه‌فقط ماشین‌های تورینگ برای محاسبهٔ
همهٔ توابع توان کافی ندارند، بلکه هیچ روش مؤثری چنین توانی ندارد.

کلید اثبات چنین نتایج سلبی این واقعیت است که می‌توانیم به خود
ماشین‌های تورینگ عدد نسبت دهیم. ساده‌ترین راه، برشمردن آن‌هاست؛ شاید
با تثبیت شیوه‌ای مشخص برای نوشتن ماشین‌های تورینگ و برنامه‌هایشان و
سپس فهرست‌کردن منظم آن‌ها. هنگامی که ببینیم این کار ممکن است، وجود
توابع تورینگ‌محاسبه‌ناپذیر از ملاحظات سادهٔ کاردینالی نتیجه می‌شود:
مجموعهٔ توابع از~$\Nat$ به~$\Nat$ (در واقع حتی فقط از~$\Nat$ به
$\{0,1\}$) !!{nonenumerable} است، اما چون همهٔ ماشین‌های تورینگ را
می‌توان برشمرد، مجموعهٔ توابع تورینگ‌محاسبه‌پذیر فقط
!!{denumerable}.

همچنین می‌توانیم توابع و مسائلی \emph{مشخص} تعریف کنیم و ثابت کنیم
که به‌ترتیب محاسبه‌ناپذیر و تصمیم‌ناپذیرند. یکی از آن‌ها \emph{مسئلهٔ
توقف} است. ماشین‌های تورینگ را می‌توان با فهرست‌کردن دستورهایشان به
طور متناهی توصیف کرد. چنین توصیفی از یک ماشین تورینگ، یعنی برنامهٔ
ماشین تورینگ، البته می‌تواند ورودیِ ماشین تورینگ دیگری باشد. پس
می‌توانیم ماشین‌هایی را در نظر بگیریم که دربارهٔ ماشین‌های دیگر
تصمیم می‌گیرند. پرسشی به‌ویژه جالب این است: «آیا ماشین تورینگ داده‌شده
پس از آغاز با ورودی~$n$ سرانجام متوقف می‌شود؟» خوب می‌بود اگر ماشین
تورینگی می‌توانست این پرسش را تصمیم کند؛ آن را همچون ماشین کنترل
کیفیتی در نظر بگیرید که مراقب است ماشین‌های تورینگ در حلقه‌های
نامتناهی و مانند آن گرفتار نشوند. واقعیت جالبی که تورینگ ثابت کرد
این است که چنین ماشینی نمی‌تواند وجود داشته باشد. هیچ ماشین تورینگ
واحدی وجود ندارد که هرگاه با ورودیِ شامل توصیف یک ماشین تورینگ~$M$ و
عددی چون~$n$ آغاز شود، همیشه متوقف شود و بسته به اینکه آیا $M$ با
ورودی~$n$ متوقف می‌شد یا نه خروجی~$1$ یا~$0$ بدهد.

هنگامی که نمونه‌هایی مشخص از مسائل تصمیم‌ناپذیر داشته باشیم، می‌توانیم
با آن‌ها تصمیم‌ناپذیری مسائل دیگری را نیز نشان دهیم. برای نمونه، یکی
از مسائل مشهورِ تصمیم‌ناپذیر این پرسش است: «آیا !!{formula}
مرتبه‌اولِ~$!A$ معتبر است؟» هیچ ماشین تورینگی وجود ندارد که با دریافت
یک !!{formula} مرتبه‌اولِ~$!A$ تضمین شود متوقف شود و بسته به معتبر بودن
یا نبودن~$!A$ خروجی~$1$ یا~$0$ بدهد. از نظر تاریخی، پرسش یافتن روشی
برای حل مؤثر این مسئله صرفاً «مسئلهٔ» تصمیم نامیده می‌شد؛ ازاین‌رو
می‌گوییم مسئلهٔ تصمیم حل‌ناپذیر است. تورینگ و چرچ این نتیجه را تقریباً
هم‌زمان و مستقل از یکدیگر ثابت کردند، و بنابراین آن را قضیهٔ
چرچ--تورینگ نیز می‌نامند.

\end{document}
